\section{Experimental Setup}
\label{sec:setup}

\subsection{Scenarios}
The agent architecture was evaluated across two main configurations based on the challenge scenarios, specifically designed to test planning efficiency and multi-agent coordination:

\begin{itemize}
    \item \textbf{Single-Agent Benchmarking:} Evaluated on the \textbf{first challenge} scenario to isolate and compare planning performance. This setup contrasts the default configuration using Breadth-First Search (BFS) against the proposed PDDL planner. 
    \item \textbf{Multi-Agent Coordination:} Evaluated on the \textbf{second challenge} scenario to analyze team dynamics. This setup compares a baseline deployment operating without communication against a cooperative team of two agents, where both utilize PDDL planning but their architectures differ: one agent operates solely via BDI, while the other is enhanced with BDI and LLM integration.
\end{itemize}

To measure the performance and efficiency of each configuration, we consider the \textbf{net score} ($R_{\text{net}} = \text{reward} - \text{penalties}$) accumulated by the agents at the end of each map, following a 5-minute run per map.

\subsection{LLM Component Evaluation Setup}
The LLM component is evaluated in isolation via a dedicated test suite (\texttt{tests/llm\_tests.ts}), with test cases stored as JSON files under \texttt{tests/data/}. Rather than evaluating the system as a monolithic block, the suite isolates individual pipeline stages across \textbf{315 total test cases} spanning 11 categories to enable a fine-grained diagnosis of failure modes.

The dataset covers every stage of the LLM pipeline. Deterministic categories include the tool handlers (\texttt{tools}, 53 cases) and the action label classifier (\texttt{extractAction}, 25 cases). LLM-backed categories target each processing step individually: math expression isolation and placeholder resolution (\texttt{planTasks}, 30 cases; \texttt{extractMathExpression}, 25 cases), request decomposition (\texttt{splitMessage}, 35 cases), geographic entity extraction (\texttt{extractCityName}, 25 cases), structured JSON constraint parsing for delivery directions (\texttt{extractDeliveryConstraint}, 25 cases) and parcel-stack rules (\texttt{extractStackConstraint}, 8 cases), tile desire generation including sub-action discrimination (\texttt{generateDesire}, 23 cases), orchestrator action classification across all eight action types (\texttt{decideNextAction}, 44 cases), and multi-agent coordination command extraction covering all four sub-commands and both explicit and implicit stop/resume signals (\texttt{extractMultiAgentCommand}, 22 cases). Ground-truth verification uses exact string matching for classification labels and specialized JSON structural validation for complex constraint objects.